\section{Scope and non-claims}\label{sec:scope}
% statements-of-record rows resolved here: none

The preceding sections build a typed apparatus for comparing the usual
BSD/Bloch--Kato scalar packages with finer source carriers.  The apparatus
is deliberately limited.  It separates what is proved at the finite typed
level from what is supplied by external arithmetic input, and it records
several audit no-gos without turning them into mathematical non-existence
claims.

The four non-claims governing this paper are the following.

\begin{itemize}
  \item \textbf{A-NC-1 --- No-gos are audit results, not non-existence claims.}
  The no-go suite (NG-1--NG-5, higher-rank, global, support-prime)
  records where theorem-grade Mode-A literature does not reach; it is not
  a claim that the missing identities cannot exist.  Here ``Mode-A''
  denotes the theorem-grade arithmetic literature audited as external
  provenance, as opposed to a finite diagnostic, numerical example, or
  support table.

  \item \textbf{A-NC-2 --- Finite typed level.}
  Adequacy collapses are proved at the finite, mathlib-free, typed level
  (identity-matrix action; \(\Q\) rank-1 projector traces);
  operator/motivic full generality is not claimed.  Thus the Schur
  residual calculations in this paper are finite vector and matrix
  computations over explicit typed carriers.  They are not asserted as
  analytic operator theorems or as constructions of motivic cycles.

  \item \textbf{A-NC-3 --- \(\vsrc\) structural narrowing.}
  The \(\vsrc\) exterior-product relations are represented structurally on
  the \(2^r\) coordinate space; the dimension/calibration content is
  faithful, full axiomatization is not claimed.  The \(\vsrc\) section
  therefore calibrates the determinant
  shape of a rank-\(r\) source and the low-rank relations used in the
  paper; it does not present a complete general exterior-algebra
  formalization or a descent to actual motivic cycles.

  \item \textbf{A-NC-4 --- Recognition-source transports are carriers.}
  The five per-column transports are recognition sources (carriers), not
  derived theorems.  They are supplied or imported as named arithmetic
  comparison content and carried as hypotheses where needed; they are not
  established from the finite typed apparatus itself.
\end{itemize}

Consequently, this paper should not be read as proving BSD, an Iwasawa main
conjecture, or a Gross--Zagier theorem.  The column computations prove
adequacy or residual statements for the typed carriers introduced here.
The recognition sources and imports are supplied, imported, or carried as
hypotheses; they are never claimed to be established by the apparatus.
